\chapter{Data types}

\section{Core}

\section{Basic}
    
    \subsection{Bytes}
        This is the rawest form of data
        \begin{minted}{python}
            raw = b'\x48\x65\x6c\x6c\x6f'
            raw.decode("utf-8")  # "Hello"
            data_bytes = b'{"name": "Ali", "age": 30}'
            data_dict = json.loads(data_bytes.decode("utf-8"))
            print(data_dict["name"])  # Ali
        \end{minted}
        
        \paragraph{As data type hint}
            \begin{minted}{python}
                class Interface(Protocol):
                    _ram: bytes
            \end{minted}
        
        \paragraph{As data type}
            \begin{minted}{python}
                data_bytes = b'hello'
                data_str = data_bytes.decode("utf-8")
                print(data_str)   # "hello"
            \end{minted}
        
        \paragraph{Converting to string}
            \begin{minted}{python}
                data_str = "hello"
                data_bytes = data_str.encode("utf-8")
                print(data_bytes)  # b'hello'
            \end{minted}
    
    \subsection{Numbers}
        
        \paragraph{int}
        
        \paragraph{float}
        
        \paragraph{complex}
        
        \paragraph{Decimal}
            for arbitary precision fro scientific calculation
        
        \paragraph{Fraction}
            For rational numbers

\section{Mutability}
    If the variable after the creation is modifiable or not

\section{Iterable}
    \begin{minted}{python}
    \end{minted}
    \subsection{Collections}
        Collection class inherits from \enquote("Iterable") so you loop  and it is also a Container so using \enquote{"in"} you can check membership and Sized so you can call len on it
        \subsubsection{Sequence}
            They have order
        
            \paragraph{list}
            
            \paragraph{str}
            \paragraph{range}
        
            \subsubsubsection{tuple}
                \paragraph{Tuple type hint}
                    Tuple hint
                    \begin{minted}{python}
                        Tuple[int,...]
                        t: Tuple[int, str, float] = (1, "a", 3.5)
                    \end{minted}
                \paragraph{named tuple}
                    \begin{minted}{python}
                        from collections import namedtuple
                        Point = namedtuple("Point", ["x", "y"])
                        p1 = Point(10, 20)
                        print(p1.x)  # 10
                        print(p1.y)  # 20
                        print(p1[0]) # 10
                    \end{minted}
                    \begin{minted}{python}
                        from typing import NamedTuple
                        class Point(NamedTuple):
                            x: float
                            y: float
                        p = Point(1.5, 2.7)
                        print(p.x, p.y)
                    \end{minted}
        \subsubsection{Shared functions}
            \paragraph{slice}
                \begin{minted}{python}
                    s = slice(2, 5, 1)
                    nums = [0, 1, 2, 3, 4, 5, 6, 7, 8, 9]
                    print(nums[s])
                \end{minted}
        \subsubsection{set types}
            \paragraph{set}
            \paragraph{fozenset}
        
        \subsubsection{Mapping types}
            Mapping type is a key value type
            \paragraph{dict}
            \paragraph{defaultdict}
                If the key doesnt exist it creates it
            \paragraph{ChainMap}
                Combines two dicts
            \paragraph{MappingProxyType}
                A read only dict
            \paragraph{Ordered Dictoionary}
                \begin{minted}{python}
                    from collections import OrderedDict
        
                    od = OrderedDict()
                    od["a"] = 1
                    od["b"] = 2
                    od["c"] = 3
                    
                    print(od)  # OrderedDict([('a', 1), ('b', 2), ('c', 3)])
        
                    
                    od.move_to_end("b")
                    print(od)  # OrderedDict([('a', 1), ('c', 3), ('b', 2)])
        
                    od.move_to_end("b", last=False)
                    print(od)  # OrderedDict([('b', 2), ('a', 1), ('c', 3)])
                \end{minted}
    

\section{Hashability}
    A hashable object has implemented \pyl{__hash__()} and \pyl{__eq__()} functions. Data types such as int, str, float are hashable but list, dict and set are not
    \paragraph{\_\_hash\_\_()}
        when implemented in a class, then it should return an  integer value sokmehow related to that object

\section{methods}
    
    remove(element): raises key error if element doesnt exist
    
    discard(element): removes but does nothing if the element doesnt exist
    
    pop: removes and returns an arbitrary element
    
    clear: removes all the elements


    